\title[Hidden attractors as a rare-event problem]{Searching for hidden attractors as a rare-event problem: An adaptive strategy portfolio with a worst-case guarantee}

\author{Maksim Elisov}
\email{maksimelisov2003@mail.ru}
\affiliation{Samara National Research University, Moskovskoye Shosse 34, Samara 443086, Russia}

\date{\today}

\begin{abstract}
An attractor is hidden if its basin does not intersect small neighborhoods of equilibria, so trajectories launched from unstable equilibria miss it; when the basin is also small, locating the attractor becomes rare-event detection in the space of initial conditions. We characterize a search method by its detection probability $P(N)$ after $N$ integrations and its cost $N_{0.9}$, relative to uniform sampling with $N_{0.9}\approx2.3/p$ for basin measure $p$. An adaptive method can gain only from what trajectories that miss the target reveal. We propose a portfolio of subset simulation driven by transient duration, active learning of basin boundaries, and a uniform share $\varepsilon$, and prove that it is at most $1/\varepsilon$ times costlier than uniform sampling while inheriting the scaling exponent of its splitting arm. On a four-level benchmark with a verification protocol, the portfolio reduces the cost of finding rare hidden attractors ($p\sim10^{-3}$--$10^{-2}$) by factors of 3--7. On a synthetic family, the cost scales as $p^{-0.45}$ down to $p=10^{-5}$, below which the time resolution of the transients limits it. On coupled copies of systems with hidden attractors, the gain in integrations reaches $10^3$--$10^4$. Measured in wall time, including against uniform sampling with recurrence-based early stopping, the gain persists for bounded trajectories (10--64 for a rare hidden regime of two coupled Chua circuits) but shrinks where trajectories escape quickly. The enumeration of regimes by the protocol reveals a hidden chaotic attractor of the coupled circuits born by the coupling. Documented failures stay within the guarantee.
\end{abstract}

\maketitle

\begin{quotation}
Some attractors of nonlinear systems are hidden: no trajectory started near an equilibrium ever reaches them, so the textbook way of finding attractors misses them. When such an attractor also has a tiny basin, finding it is like searching for a needle in a haystack of initial conditions, and published methods rarely report whether they beat random guessing. We recast the search as a rare-event problem and measure the cost of finding a target with a prescribed confidence. This makes methods comparable and shows that an advantage can come only from what unsuccessful trajectories reveal about the basins. On this basis, we build a search strategy that is guaranteed never to be much worse than random sampling, is several times better on typical problems, and is orders of magnitude better on rare targets in coupled systems, together with a benchmark and a protocol that catch numerical artifacts and transient chaos mistaken for attractors. Measured in computing time rather than in the number of trajectories, the advantage holds where trajectories stay bounded, as in typical physical systems.
\end{quotation}
